% Generated semantic source for AI Almanach v3.
% Edit v3 directly after initial generation; do not overwrite
% editorial changes by rerunning the converter casually.

% ==== Programming Foundations =============================================================================
\chapter{Programming Foundations}
\chapterlead{How does an analytical idea become reproducible, testable software?}{Specify}{Implement}{Test}{Maintain}{programming}

\label{ch:programming-foundations}

\begin{learningobjectives}
After working through this chapter you should be able to:
\begin{itemize}
  \item predict the result shape of a broadcast operation before running it,
        and recognise the silent bug that broadcasting causes most often;
  \item explain why \texttt{0.1 + 0.2 == 0.3} is \texttt{False} and what to
        write instead;
  \item name the three Python defaults that produce wrong results rather than
        errors;
  \item choose between a loop, a vectorised operation, and a compiled routine
        from the data size;
  \item write the environment record that makes an experiment reproducible;
  \item read a shape error and locate the bug without a debugger.
\end{itemize}
\end{learningobjectives}

\begin{plainlanguage}
The programming you need for machine learning is narrower than general software
engineering but less forgiving in one specific way: most mistakes do not crash.
A wrong axis, a broadcast that silently expands, a copy that was actually a
view, a seed that was never set --- each produces a number, and the number
looks fine. This chapter is mostly about the failures that return a result
instead of an error.
\end{plainlanguage}

\begin{engineernote}
Two habits prevent most of the damage. First, print shapes --- of every array,
at every boundary, until the pipeline is trusted. Second, write the assertion
you are tempted to check by eye: \texttt{assert y\_pred.shape == y\_true.shape}
costs one line and catches a class of bug that can otherwise survive all the
way into a published number.
\end{engineernote}

\section{Programming with Python}
\subsection{Introduction to Python}

% ---- 3.1.1 Data structures ----------------------------------------------------------------------
\subsubsection{Data structures}

\begin{v3topic}{Data Structures: Lists, Tuples, Dicts, Sets}
    Python provides four built-in collection types that differ in mutability, ordering, and uniqueness\textsuperscript{\cite{romanoLearnPythonProgramming2021}}.
    Choosing the right type matters for both correctness and performance.
\visualseparator
    \begin{tblr}{width=\linewidth,colspec={|Q[l,m,2cm]|Q[c,m]|Q[c,m]|X|}}
        \hline
        \textbf{Type} & \textbf{Ordered} & \textbf{Mutable} & \textbf{Notes / Use-case} \\
        \hline
        \texttt{list}  & yes & yes & General-purpose sequence; \texttt{[1, 2, 3]} \\
        \hline
        \texttt{tuple} & yes & no  & Immutable record; unpacking; function return values \\
        \hline
        \texttt{dict}  & yes (3.7+) & yes & Key $\to$ value mapping; $O(1)$ lookup \\
        \hline
        \texttt{set}   & no  & yes & Unique elements; fast membership tests; set algebra \\
        \hline
    \end{tblr}
\captionof{table}{The four built-in collections. Choose by the two middle columns: ordering and mutability decide correctness, everything else is performance.}
\label{tab:prog-data-structures}
\end{v3topic}

% ---- 3.1.2 Functions ----------------------------------------------------------------------------
\subsubsection{Functions}

\begin{v3topic}{Defining Functions}
    Functions are first-class objects in Python\textsuperscript{\cite{romanoLearnPythonProgramming2021}}.
    They are defined with \texttt{def}, accept positional and keyword arguments, and always return a value (\texttt{None} if no \texttt{return} statement).
\visualseparator
    \begin{itemize}
        \item \texttt{def f(x, y=0):} --- positional + default argument
        \item \texttt{*args} --- collects extra positionals as a tuple
        \item \texttt{**kwargs} --- collects extra keywords as a dict
        \item \texttt{lambda x: x**2} --- anonymous single-expression function
        \item Multiple return values via tuple: \texttt{return a, b}
    \end{itemize}

\end{v3topic}
\begin{v3topic}{Higher-Order Functions}
    Because functions are objects, they can be passed as arguments, returned from other functions, and stored in data structures.
    The built-ins \texttt{map()}, \texttt{filter()}, and \texttt{sorted(key=...)} are the most common higher-order functions in practice\textsuperscript{\cite{romanoLearnPythonProgramming2021}}.
\visualseparator
    \begin{itemize}
        \item \texttt{map(fn, iterable)} --- apply \texttt{fn} to each element
        \item \texttt{filter(pred, iterable)} --- keep elements where \texttt{pred} is true
        \item \texttt{sorted(lst, key=lambda x: x[1])} --- sort by key
        \item Closures: inner function captures enclosing variable
        \item \texttt{functools.partial} --- fix arguments of a function
    \end{itemize}
\end{v3topic}

% ---- 3.1.3 Flow control -------------------------------------------------------------------------
\subsubsection{Flow control}

\begin{v3topic}{Flow Control and Comprehensions}
    Python's control flow is concise: \texttt{if}/\texttt{elif}/\texttt{else} for branching, \texttt{for} over any iterable, \texttt{while} for condition-driven loops\textsuperscript{\cite{romanoLearnPythonProgramming2021}}.
    List, dict, and set comprehensions replace many explicit loops with readable one-liners.
\visualseparator
    \begin{tblr}{width=\linewidth,colspec={|X[2]|X[3]|}}
        \hline
        \textbf{Construct} & \textbf{Example} \\
        \hline
        Conditional      & \texttt{x = "yes" if n > 0 else "no"} \\
        \hline
        for-loop         & \texttt{for i, v in enumerate(lst):} \\
        \hline
        while-loop       & \texttt{while q: q.pop()} \\
        \hline
        List comp.       & \texttt{[x**2 for x in range(10) if x\%2==0]} \\
        \hline
        Dict comp.       & \texttt{\{k: v for k, v in pairs\}} \\
        \hline
        Generator expr.  & \texttt{sum(x**2 for x in data)} --- lazy evaluation \\
        \hline
    \end{tblr}
\captionof{table}{Flow control constructs. Comprehensions are preferred not for brevity but because they make the output shape explicit at the point of construction.}
\label{tab:prog-flow-control}
\end{v3topic}

% ---- 3.1.4 Input / Output -----------------------------------------------------------------------
\subsubsection{Input / Output}

\begin{v3topic}{File Handling and I/O}
    Python's \texttt{open()} built-in and the \texttt{with} context manager handle file I/O safely: the file is closed automatically when the block exits, even on exceptions\textsuperscript{\cite{romanoLearnPythonProgramming2021}}.
\visualseparator
    \begin{tblr}{width=\linewidth,colspec={|Q[l,m,2.5cm]|X|}}
        \hline
        \textbf{Pattern} & \textbf{Description} \\
        \hline
        \texttt{with open(path, "r") as f:} & Read text file; auto-close on exit \\
        \hline
        \texttt{f.read()} / \texttt{f.readlines()} & Read entire content or list of lines \\
        \hline
        \texttt{open(path, "w")} & Write (truncates); \texttt{"a"} to append \\
        \hline
        \texttt{open(path, "rb")} & Binary mode for images, serialised objects, etc. \\
        \hline
        \texttt{print(x, file=sys.stderr)} & Write to standard error stream \\
        \hline
        \texttt{json.dump} / \texttt{json.load} & Serialize Python objects to/from JSON \\
        \hline
    \end{tblr}
\captionof{table}{File-handling patterns. The context-manager form is the only one that closes the handle when an exception is raised part-way through.}
\label{tab:prog-file-io}
\end{v3topic}

% ---- 3.1.5 Modules and packages -----------------------------------------------------------------
\subsubsection{Modules and packages}

\begin{v3topic}{Modules and the Import System}
    A \emph{module} is a \texttt{.py} file; a \emph{package} is a directory containing an \texttt{\_\_init\_\_.py} file.
    Python resolves imports by searching \texttt{sys.path}: the script's directory, environment packages, and standard library\textsuperscript{\cite{romanoLearnPythonProgramming2021}}.
\visualseparator
    \begin{itemize}
        \item \texttt{import numpy as np} --- import with alias
        \item \texttt{from pathlib import Path} --- selective import
        \item \texttt{\_\_init\_\_.py} --- marks directory as package; runs on import
        \item \texttt{if \_\_name\_\_ == "\_\_main\_\_":} --- guard for runnable modules
        \item Circular imports: break by restructuring or lazy imports
    \end{itemize}

\end{v3topic}
\begin{v3topic}{Package Management with pip}
    Third-party packages are installed from PyPI using \texttt{pip}\textsuperscript{\cite{romanoLearnPythonProgramming2021}}.
    Reproducible environments are specified via \texttt{requirements.txt} or the modern \texttt{pyproject.toml}.
\visualseparator
    \begin{itemize}
        \item \texttt{pip install numpy==1.26.0} --- install pinned version
        \item \texttt{pip freeze > requirements.txt} --- capture current environment
        \item \texttt{pip install -r requirements.txt} --- recreate environment
        \item \texttt{pyproject.toml} --- modern standard (PEP 517/518) for project metadata and deps
    \end{itemize}
\end{v3topic}

% ---- 3.2 Classes and Inheritance ----------------------------------------------------------------
\subsection{Classes and Inheritance}

% ---- 3.2.1 Scopes and namespaces ----------------------------------------------------------------
\subsubsection{Scopes and namespaces}

\begin{v3topic}{The LEGB Rule}
    Python resolves names by searching scopes in order: \textbf{L}ocal $\to$ \textbf{E}nclosing $\to$ \textbf{G}lobal $\to$ \textbf{B}uilt-in (LEGB)\textsuperscript{\cite{romanoLearnPythonProgramming2021}}.
    Assignment in a scope creates a local binding; \texttt{global} and \texttt{nonlocal} keywords break out of the default rule.
\visualseparator
    \begin{tblr}{width=\linewidth,colspec={|Q[l,m,2cm]|X|}}
        \hline
        \textbf{Scope} & \textbf{Where it lives} \\
        \hline
        Local (L)     & Inside the current function call \\
        \hline
        Enclosing (E) & Surrounding function(s) --- relevant for closures \\
        \hline
        Global (G)    & Module-level namespace (\texttt{global x} to write) \\
        \hline
        Built-in (B)  & Python's built-in names (\texttt{len}, \texttt{range}, \ldots) \\
        \hline
    \end{tblr}
\captionof{table}{Python's name resolution order. Nearly every surprising scope bug is a name found at an outer level when an inner one was intended.}
\label{tab:prog-legb}
\end{v3topic}

% ---- 3.2.2 Classes and inheritance --------------------------------------------------------------
\subsubsection{Classes and inheritance}

\begin{v3topic}{Classes and Inheritance}
    Python uses \emph{class-based} OOP.
    \texttt{\_\_init\_\_} is the constructor; \texttt{self} is the explicit reference to the instance.
    Single and multiple inheritance are supported; method resolution follows the \emph{C3 linearization} (MRO)\textsuperscript{\cite{romanoLearnPythonProgramming2021}}.
\visualseparator
    \begin{tblr}{width=\linewidth,colspec={|Q[l,m,3.5cm]|X|}}
        \hline
        \textbf{Concept} & \textbf{Description} \\
        \hline
        \texttt{class Dog(Animal):}         & Inherit from \texttt{Animal} \\
        \hline
        \texttt{def \_\_init\_\_(self, ...):} & Constructor; initialise instance attributes \\
        \hline
        \texttt{super().\_\_init\_\_(...)}    & Delegate to parent constructor \\
        \hline
        \texttt{Dog.\_\_mro\_\_}              & Tuple of classes in method resolution order \\
        \hline
        \texttt{@classmethod} / \texttt{@staticmethod} & Alternate constructor or utility method \\
        \hline
        \texttt{@property}                  & Getter/setter without explicit call parentheses \\
        \hline
    \end{tblr}
\captionof{table}{Core object-oriented concepts. In machine-learning code these matter mainly for dataset and model classes, where framework base classes fix the interface.}
\label{tab:prog-classes}
\end{v3topic}

% ---- 3.2.3 Iterators and generators -------------------------------------------------------------
\subsubsection{Iterators and generators}

\begin{v3topic}{Iterators and Generators}
    Any object implementing \texttt{\_\_iter\_\_} and \texttt{\_\_next\_\_} is an \emph{iterator}.
    A \emph{generator function} uses \texttt{yield} to produce values lazily, suspending execution between calls\textsuperscript{\cite{romanoLearnPythonProgramming2021}}.
    Generators avoid loading entire sequences into memory --- essential for large datasets.
\visualseparator
    \begin{itemize}
        \item \texttt{iter(obj)} / \texttt{next(obj)} --- protocol entry points
        \item Generator function: \texttt{def gen(): yield x}
        \item Generator expression: \texttt{(x**2 for x in data)}
        \item \texttt{itertools} module: \texttt{chain}, \texttt{islice}, \texttt{product}, \ldots
    \end{itemize}

\end{v3topic}
\begin{v3topic}{Protocol Dunder Methods}
    Python's data model exposes behaviour through special (``dunder'') methods\textsuperscript{\cite{romanoLearnPythonProgramming2021}}.
    Implementing them makes custom objects behave like built-in types.
\visualseparator
    \noindent\begin{tblr}{width=\linewidth,colspec={Q[l,m,3cm]X},
                           hlines,vlines}
        \hline
        \textbf{Method} & \textbf{Triggered by} \\
        \hline
        \texttt{\_\_len\_\_}    & \texttt{len(obj)} \\
        \hline
        \texttt{\_\_getitem\_\_} & \texttt{obj[key]}, slicing, iteration \\
        \hline
        \texttt{\_\_contains\_\_} & \texttt{x in obj} \\
        \hline
        \texttt{\_\_repr\_\_} / \texttt{\_\_str\_\_} & \texttt{repr(obj)} / \texttt{str(obj)} \\
        \hline
        \texttt{\_\_enter\_\_} / \texttt{\_\_exit\_\_} & \texttt{with} statement \\
        \hline
    \end{tblr}
\captionof{table}{Dunder methods and the syntax that triggers them. Implementing these is how a custom class becomes usable by code that never heard of it.}
\label{tab:prog-dunder}
\end{v3topic}

% ---- 3.3 Errors and Exceptions ------------------------------------------------------------------
\subsection{Errors and Exceptions}

\begin{workedexample}[Three Python results that are wrong rather than broken]
\emph{1. Floating-point equality.}
\[
  0.1 + 0.2 = 0.30000000000000004 \neq 0.3 .
\]
Neither $0.1$ nor $0.2$ is representable in binary, so the sum lands one unit
in the last place away from $0.3$ (see \cref{sec:math-numerics}). Never compare
floats with \texttt{==}; compare with a tolerance, e.g.\
\texttt{math.isclose(a, b, rel\_tol=1e-9)}.

\emph{2. Mutable default arguments.} A default is evaluated \emph{once}, when
the function is defined --- not on each call:
\begin{center}
\texttt{def add(item, bucket=[]): bucket.append(item); return bucket}
\end{center}
Calling \texttt{add(1)} returns \texttt{[1]}; calling \texttt{add(2)}
returns \texttt{[1, 2]}. The list persists across calls. Use
\texttt{bucket=None} and create the list inside.

\emph{3. Integer division and silent type change.} \texttt{7 / 2} is
\texttt{3.5} but \texttt{7 // 2} is \texttt{3}, and mixing the two inside an
index computation shifts every subsequent result by one without complaint.

All three return a value. None of them raises. That is what makes them
expensive.
\end{workedexample}

% ---- 3.3.1 Syntax errors ------------------------------------------------------------------------
\subsubsection{Syntax errors}

\begin{v3topic}{Syntax vs.\ Runtime Errors}
    Python distinguishes errors that prevent parsing (\emph{SyntaxError}) from errors raised during execution (\emph{exceptions})\textsuperscript{\cite{romanoLearnPythonProgramming2021}}.
\visualseparator
    \noindent\begin{tblr}{width=\linewidth,
                           colspec={Q[l,m,3.2cm]X},hlines,vlines}
        \hline
        \textbf{Error type} & \textbf{Cause} \\
        \hline
        \texttt{SyntaxError}    & Invalid Python grammar, e.g.\ missing colon, mismatched bracket \\
        \hline
        \texttt{IndentationError} & Wrong indentation level (subset of \texttt{SyntaxError}) \\
        \hline
        \texttt{NameError}      & Using a name before it is defined \\
        \hline
        \texttt{TypeError}      & Operation on incompatible types \\
        \hline
        \texttt{ValueError}     & Right type, wrong value, e.g.\ \texttt{int("abc")} \\
        \hline
        \texttt{IndexError} / \texttt{KeyError} & Access beyond sequence / absent dict key \\
        \hline
    \end{tblr}
\captionof{table}{Errors by when they surface. Syntax errors are free; runtime errors cost a run; the dangerous third category, wrong results without an error, appears in neither row.}
\label{tab:prog-error-types}
\end{v3topic}

% ---- 3.3.2 Handling and raising exceptions ------------------------------------------------------
\subsubsection{Handling and raising exceptions}

\begin{v3topic}{try / except / finally}
    The \texttt{try} block wraps code that may fail; \texttt{except} catches specific exceptions; \texttt{else} runs when no exception occurred; \texttt{finally} always runs (clean-up)\textsuperscript{\cite{romanoLearnPythonProgramming2021}}.
\visualseparator
    \begin{itemize}
        \item \texttt{except ValueError as e:} --- bind exception to name
        \item \texttt{except (TypeError, KeyError):} --- catch multiple types
        \item \texttt{raise} --- re-raise current exception
        \item \texttt{raise ValueError("bad input")} --- raise with message
        \item \texttt{finally:} --- cleanup regardless of outcome (close files, release locks)
    \end{itemize}

\end{v3topic}
\begin{v3topic}{Context Managers as Error Guards}
    The \texttt{with} statement uses \texttt{\_\_enter\_\_} / \texttt{\_\_exit\_\_} to guarantee resource release, even when exceptions occur\textsuperscript{\cite{romanoLearnPythonProgramming2021}}.
    \texttt{contextlib.contextmanager} lets you write a generator-based context manager.
\visualseparator
    \begin{itemize}
        \item \texttt{with open(f) as fh:} --- file closed on any exit
        \item \texttt{with contextlib.suppress(OSError):} --- silently skip errors
        \item \texttt{@contextlib.contextmanager} decorator: \texttt{yield} separates enter from exit logic
    \end{itemize}
\end{v3topic}

% ---- 3.3.3 User-defined exceptions --------------------------------------------------------------
\subsubsection{User-defined exceptions}

\begin{v3topic}{Custom Exception Classes}
    User-defined exceptions inherit from \texttt{Exception} (or a more specific base class).
    Custom exceptions improve code clarity and allow callers to catch domain-specific errors without catching everything\textsuperscript{\cite{romanoLearnPythonProgramming2021}}.
\visualseparator
    \begin{tblr}{width=\linewidth,colspec={|Q[l,m,4cm]|X|}}
        \hline
        \textbf{Pattern} & \textbf{Notes} \\
        \hline
        \texttt{class DataError(Exception): pass}       & Minimal custom exception \\
        \hline
        \texttt{class DataError(ValueError): ...}       & Specialise a standard exception \\
        \hline
        Add \texttt{\_\_init\_\_} to store attributes   & e.g.\ \texttt{self.column = column} for context \\
        \hline
        Application hierarchy                           & \texttt{AppError} $\to$ \texttt{DBError}, \texttt{NetworkError} \\
        \hline
    \end{tblr}
\captionof{table}{Custom exception patterns. A specific exception type lets a caller handle one failure and let the rest propagate, which a bare string message cannot.}
\label{tab:prog-custom-exceptions}
    Catch the base class to handle the whole hierarchy; catch subclasses for fine-grained handling.
\end{v3topic}

% ---- 3.4 Important Libraries --------------------------------------------------------------------
\subsection{Important Libraries}

% ---- 3.4.1 Standard Python library --------------------------------------------------------------
\subsubsection{Standard Python library}

\begin{v3topic}{Python Standard Library Highlights}
    The standard library ships with CPython and covers a huge range of tasks without third-party dependencies\textsuperscript{\cite{romanoLearnPythonProgramming2021}}.
\visualseparator
    \begin{tblr}{width=\linewidth,colspec={|Q[l,m,2.5cm]|X|}}
        \hline
        \textbf{Module} & \textbf{Key functionality} \\
        \hline
        \texttt{os} / \texttt{os.path} & OS interface: env vars, paths, processes \\
        \hline
        \texttt{pathlib}     & Object-oriented path handling (\texttt{Path("data") / "file.csv"}) \\
        \hline
        \texttt{sys}         & Interpreter state: \texttt{sys.argv}, \texttt{sys.path}, \texttt{sys.exit()} \\
        \hline
        \texttt{json}        & Serialize/deserialize JSON data \\
        \hline
        \texttt{datetime}    & Date, time, timedelta arithmetic \\
        \hline
        \texttt{collections} & \texttt{Counter}, \texttt{defaultdict}, \texttt{namedtuple}, \texttt{deque} \\
        \hline
        \texttt{itertools}   & Efficient iterators: \texttt{chain}, \texttt{product}, \texttt{groupby} \\
        \hline
        \texttt{re}          & Regular expressions \\
        \hline
    \end{tblr}
\captionof{table}{Standard-library modules worth knowing before reaching for a dependency.}
\label{tab:prog-stdlib}
\end{v3topic}

% ---- 3.4.2 Scientific calculations --------------------------------------------------------------
\subsubsection{Scientific calculations}

\begin{v3topic}{NumPy and SciPy}
    \textbf{NumPy} provides the \texttt{ndarray}: a fast, homogeneously-typed N-dimensional array with vectorised operations and broadcasting\textsuperscript{\cite{fuhrerScientificComputingPython2021}}.
    \textbf{SciPy} builds on NumPy and adds algorithms for linear algebra, optimization, integration, signal processing, and statistics.
\visualseparator
    \begin{itemize}
        \item \texttt{numpy.linalg}: matrix factorisation, eigenvalues, solve
        \item \texttt{scipy.optimize}: minimize, root-finding, curve fitting
        \item \texttt{scipy.stats}: probability distributions, hypothesis tests
        \item \texttt{scipy.sparse}: sparse matrix formats and solvers
        \item \texttt{scipy.signal}: convolution, filtering, spectral analysis
    \end{itemize}

\end{v3topic}
\begin{v3topic}{Performance: Memory Layout and ufuncs}
    NumPy arrays are stored in contiguous memory (C-order by default), enabling SIMD vectorisation\textsuperscript{\cite{fuhrerScientificComputingPython2021}}.
    \emph{Universal functions} (\texttt{ufuncs}) apply element-wise operations over arrays without Python loops.
\visualseparator
    \begin{itemize}
        \item \texttt{np.dot(A, B)} --- BLAS-backed matrix multiplication
        \item \texttt{np.sum}, \texttt{np.mean}, \texttt{np.std} --- reduction ufuncs
        \item \texttt{dtype} selection: \texttt{float32} vs.\ \texttt{float64} trades precision for memory
        \item \texttt{np.einsum("ij,jk->ik", A, B)} --- Einstein summation notation
    \end{itemize}
\end{v3topic}

% ---- 3.4.3 Speeding up Python -------------------------------------------------------------------
\subsubsection{Speeding up Python}

\begin{v3topic}{Accelerating Python Code}
    Pure Python loops are slow due to interpreter overhead.
    Several tools compile or parallelize Python to approach C/Fortran speeds\textsuperscript{\cite{fuhrerScientificComputingPython2021}}.
\visualseparator
    \begin{tblr}{width=\linewidth,colspec={|Q[l,m,2.5cm]|X|}}
        \hline
        \textbf{Tool} & \textbf{Approach and notes} \\
        \hline
        \textbf{NumPy}          & Vectorise: replace Python loops with array operations \\
        \hline
        \textbf{Numba}          & JIT-compile with \texttt{@numba.jit}; no code rewrite required \\
        \hline
        \textbf{Cython}         & Annotate Python with C types; compiled to C extension \\
        \hline
        \textbf{multiprocessing} & True parallelism bypassing the GIL; \texttt{Pool.map()} \\
        \hline
        \textbf{concurrent.futures} & \texttt{ProcessPoolExecutor} for CPU; \texttt{ThreadPoolExecutor} for I/O \\
        \hline
        \textbf{Dask}           & Lazy parallel arrays/dataframes; scales to cluster \\
        \hline
    \end{tblr}
\captionof{table}{Ways to make Python faster, in increasing order of commitment. Measure before choosing: the answer is usually vectorisation, not a rewrite.}
\label{tab:prog-acceleration}
    Rule of thumb: profile first (\texttt{cProfile}, \texttt{line\_profiler}), then optimise the bottleneck.
\end{v3topic}

% ---- 3.4.4 Visualization ------------------------------------------------------------------------
\subsubsection{Visualization}

\begin{v3topic}{Python Visualization Ecosystem}
    Python offers a layered visualization stack from low-level to high-level APIs\textsuperscript{\cite{navlaniPythonDataAnalysis2021}}.
\visualseparator
    \begin{tblr}{width=\linewidth,colspec={|Q[l,m,2.5cm]|X|}}
        \hline
        \textbf{Library} & \textbf{Description} \\
        \hline
        \textbf{Matplotlib} & Foundation: fine-grained control; \texttt{fig, ax = plt.subplots()} \\
        \hline
        \textbf{Seaborn}    & Statistical plots on top of Matplotlib; tidy-data API \\
        \hline
        \textbf{Plotly}     & Interactive HTML charts; integrates with Dash for dashboards \\
        \hline
        \textbf{Bokeh}      & Browser-based interactive plots; streaming data support \\
        \hline
        \textbf{Altair}     & Declarative API based on Vega-Lite; grammar of graphics \\
        \hline
    \end{tblr}
\captionof{table}{The plotting libraries and what each is actually for. See \cref{ch:visualization} for choosing the chart rather than the library.}
\label{tab:prog-viz-libraries}
    Matplotlib is the default for publication figures; Plotly/Bokeh/Altair for interactive exploration.
\end{v3topic}

% ---- 3.4.5 Accessing databases ------------------------------------------------------------------
\subsubsection{Accessing databases}

\begin{v3topic}{Database Access in Python}
    Python accesses databases via the DB-API 2.0 standard (\texttt{sqlite3}, \texttt{psycopg2}, \texttt{pymysql}) or the ORM layer SQLAlchemy\textsuperscript{\cite{navlaniPythonDataAnalysis2021}}.
\visualseparator
    \begin{tblr}{width=\linewidth,colspec={|Q[l,m,3cm]|X|}}
        \hline
        \textbf{Tool} & \textbf{Notes} \\
        \hline
        \texttt{sqlite3}        & Stdlib; file-based SQL; good for prototyping \\
        \hline
        \texttt{psycopg2}       & PostgreSQL adapter; DB-API 2.0 compliant \\
        \hline
        \textbf{SQLAlchemy Core} & SQL expression language; database-agnostic \\
        \hline
        \textbf{SQLAlchemy ORM}  & Map Python classes to tables; session-based transactions \\
        \hline
        \textbf{Pandas}          & \texttt{pd.read\_sql(query, con)} --- query result directly to DataFrame \\
        \hline
    \end{tblr}
\captionof{table}{Database access options, from raw driver to full ORM. The trade is control against boilerplate.}
\label{tab:prog-database}
\end{v3topic}

% ---- 3.5 Working with Python --------------------------------------------------------------------
\subsection{Working with Python}

% ---- 3.5.1 Virtual environments -----------------------------------------------------------------
\subsubsection{Virtual environments}

\begin{v3topic}{Virtual Environments}
    A virtual environment is an isolated Python installation that keeps project dependencies separate, avoiding conflicts between projects\textsuperscript{\cite{romanoLearnPythonProgramming2021}}.
\visualseparator
    \begin{tblr}{width=\linewidth,colspec={|Q[l,m,3cm]|X|}}
        \hline
        \textbf{Tool} & \textbf{Commands / Notes} \\
        \hline
        \texttt{venv} (stdlib)   & \texttt{python -m venv .venv}; activate with \texttt{source .venv/bin/activate} \\
        \hline
        \texttt{conda}           & \texttt{conda create -n myenv python=3.11}; manages non-Python deps too \\
        \hline
        \texttt{poetry}          & Dependency resolution + virtual env + publishing in one tool \\
        \hline
        \texttt{uv}              & Ultra-fast Rust-based pip/venv replacement \\
        \hline
    \end{tblr}
\captionof{table}{Environment tools. Any of them works; using none of them is what makes a result unreproducible.}
\label{tab:prog-venv}
    Best practice: one virtual environment per project; never install into the system Python.
\end{v3topic}

% ---- 3.5.2 Managing packages --------------------------------------------------------------------
\subsubsection{Managing packages}

\begin{v3topic}{Reproducible Package Management}
    Pinning dependency versions ensures reproducible builds across machines and CI/CD pipelines\textsuperscript{\cite{romanoLearnPythonProgramming2021}}.
\visualseparator
    \begin{tblr}{width=\linewidth,colspec={|Q[l,m,3.5cm]|X|}}
        \hline
        \textbf{File / Tool} & \textbf{Purpose} \\
        \hline
        \texttt{requirements.txt}   & Simple pinned list; \texttt{pip freeze > requirements.txt} \\
        \hline
        \texttt{pyproject.toml}     & Modern standard (PEP 517/518/621); replaces \texttt{setup.py} \\
        \hline
        \texttt{pip-compile}        & Lock transitive deps to \texttt{requirements.txt} from \texttt{.in} spec \\
        \hline
        \texttt{conda env export}   & Export conda environment including non-Python packages \\
        \hline
    \end{tblr}
\captionof{table}{The files that pin an environment. A lock file records what was actually installed; a requirements file records what was requested, and the two diverge.}
\label{tab:prog-packaging}
\end{v3topic}

% ---- 3.5.3 Unit and integration testing ---------------------------------------------------------
\subsubsection{Unit and integration testing}

\begin{v3topic}{Testing with pytest}
    \texttt{pytest} is the de-facto Python testing framework: plain \texttt{assert} statements, powerful fixtures, parametrize, and a rich plugin ecosystem\textsuperscript{\cite{anayaCleanCodePython2021}}.
\visualseparator
    \begin{itemize}
        \item Test files named \texttt{test\_*.py}; functions named \texttt{test\_*}
        \item \texttt{@pytest.fixture} --- reusable setup/teardown
        \item \texttt{@pytest.mark.parametrize} --- run one test with multiple inputs
        \item \texttt{pytest.raises(ValueError)} --- assert exception is raised
        \item \texttt{pytest --cov} (pytest-cov) --- measure code coverage
    \end{itemize}

\end{v3topic}
\begin{v3topic}{Test Design Principles}
    Good tests are fast, isolated, repeatable, and self-validating\textsuperscript{\cite{anayaCleanCodePython2021}}.
    The \emph{Arrange-Act-Assert} (AAA) pattern structures each test clearly.
\visualseparator
    \begin{tblr}{width=\linewidth,colspec={|Q[l,m,2cm]|X|}}
        \hline
        \textbf{Layer} & \textbf{Scope} \\
        \hline
        Unit test       & Single function / class; mock external deps \\
        \hline
        Integration     & Several components together; real DB or filesystem \\
        \hline
        End-to-end      & Full system from API to DB and back \\
        \hline
        Property test   & \texttt{hypothesis} library generates random inputs \\
        \hline
    \end{tblr}
\captionof{table}{The test pyramid. Fast, numerous unit tests catch most regressions; slow end-to-end tests catch the ones that matter most.}
\label{tab:prog-test-layers}
\end{v3topic}

% ---- 3.5.4 Documenting code ---------------------------------------------------------------------
\subsubsection{Documenting code}

\begin{v3topic}{Docstrings and Type Hints}
    Python supports inline documentation through \emph{docstrings} (PEP 257) and static type annotations (PEP 484)\textsuperscript{\cite{anayaCleanCodePython2021}}.
    Tools such as \texttt{Sphinx} generate HTML/PDF API docs from docstrings; \texttt{mypy} and \texttt{pyright} perform static type checking.
\visualseparator
    \begin{tblr}{width=\linewidth,colspec={|Q[l,m,3cm]|X|}}
        \hline
        \textbf{Convention} & \textbf{Notes} \\
        \hline
        Google style        & Sections: Args, Returns, Raises; readable prose \\
        \hline
        NumPy style         & Parameters / Returns with type + description lines \\
        \hline
        \texttt{def f(x: int) -> float:} & Type hint syntax; checked by \texttt{mypy} \\
        \hline
        \texttt{dataclasses.\allowbreak dataclass} & Auto-generates
        \texttt{\_\_init\_\_}, \texttt{\_\_repr\_\_} from annotated fields \\
        \hline
        \texttt{sphinx-autodoc}          & Extracts docstrings to build HTML reference docs \\
        \hline
    \end{tblr}
\captionof{table}{Documentation and typing conventions. Type hints do not execute, but they let a checker find shape and type errors before a run costs GPU time.}
\label{tab:prog-docstrings}
\end{v3topic}

% ---- 3.6 Version Control ------------------------------------------------------------------------
\subsection{Version Control}

% ---- 3.6.1 Introduction to version control ------------------------------------------------------
\subsubsection{Introduction to version control}

\begin{v3topic}{Why Version Control?}
    A \emph{version control system} (VCS) records changes to files over time, enabling rollback, parallel development, and collaboration\textsuperscript{\cite{romanoLearnPythonProgramming2021}}.
\visualseparator
    \begin{tblr}{width=\linewidth,colspec={|Q[l,m,3cm]|X[2]|X[2]|}}
        \hline
        \textbf{Property} & \textbf{Centralized (SVN)} & \textbf{Distributed (Git)} \\
        \hline
        Repository      & Single server copy       & Full history on every machine \\
        \hline
        Offline work    & Limited                  & Full history available offline \\
        \hline
        Branching       & Expensive                & Cheap (just a pointer) \\
        \hline
        Fault tolerance & Single point of failure  & Any clone can restore the repo \\
        \hline
    \end{tblr}
\captionof{table}{Centralised against distributed version control. The practical consequence of the distributed model is that history and branching are local and therefore cheap.}
\label{tab:prog-vcs-models}
\end{v3topic}

% ---- 3.6.2 Version control with GIT -------------------------------------------------------------
\subsubsection{Version control with GIT}

\begin{v3topic}{Core Git Concepts}
    Git stores snapshots (not diffs) as a directed acyclic graph (DAG) of \emph{commits}\textsuperscript{\cite{romanoLearnPythonProgramming2021}}.
    A \emph{branch} is a movable pointer to a commit; \emph{HEAD} is the currently checked-out commit.
\visualseparator
    \begin{tblr}{width=\linewidth,colspec={|Q[l,m,2.5cm]|X|}}
        \hline
        \textbf{Command} & \textbf{Effect} \\
        \hline
        \texttt{git init}          & Create empty repository \\
        \hline
        \texttt{git add -p}        & Stage changes interactively \\
        \hline
        \texttt{git commit -m "..."} & Record snapshot with message \\
        \hline
        \texttt{git branch feat}   & Create branch pointer \\
        \hline
        \texttt{git merge feat}    & Integrate branch into current \\
        \hline
        \texttt{git rebase main}   & Replay commits on top of main \\
        \hline
    \end{tblr}
\captionof{table}{The Git commands that cover almost all daily use.}
\label{tab:prog-git-core}

\end{v3topic}
\begin{v3topic}{Collaboration Workflow}
    Remote repositories on GitHub/GitLab enable team collaboration.
    The \emph{feature-branch workflow} is the most common: work on a branch, open a pull request, review, then merge\textsuperscript{\cite{romanoLearnPythonProgramming2021}}.
\visualseparator
    \begin{tblr}{width=\linewidth,colspec={|Q[l,m,3cm]|X|}}
        \hline
        \textbf{Command} & \textbf{Effect} \\
        \hline
        \texttt{git remote add origin URL} & Link to remote \\
        \hline
        \texttt{git push -u origin feat}   & Push branch, set upstream \\
        \hline
        \texttt{git pull --rebase}         & Fetch + rebase local commits \\
        \hline
        \texttt{git stash}                 & Temporarily shelve dirty state \\
        \hline
        \texttt{git log --oneline --graph} & Visual commit DAG \\
        \hline
    \end{tblr}
\captionof{table}{Collaboration commands. Note which of these rewrite history: those are the ones that need agreement before use on a shared branch.}
\label{tab:prog-git-collab}
\end{v3topic}

% ==================================================================================================
\section{NumPy and Pandas}
\subsection{Introduction to NumPy}

\begin{workedexample}[Broadcasting, and the bug it hides]
NumPy aligns shapes from the \emph{right} and stretches any dimension of size
$1$.

\emph{The intended behaviour.}
\[
  \underbrace{(3,1)}_{\text{column}} \;+\; \underbrace{(1,4)}_{\text{row}}
  \;\longrightarrow\; (3,4).
\]
With column $[1,2,3]^{\!\top}$ and row $[10,20,30,40]$ the result is the full
outer sum
$\left(\begin{smallmatrix}11&21&31&41\\12&22&32&42\\13&23&33&43
\end{smallmatrix}\right)$ --- useful, and exactly what was asked for.

\emph{The bug.} Now suppose targets have shape $(100,)$ and predictions come
back from a model with shape $(100,1)$ --- a single trailing dimension nobody
removed. Align from the right:
\[
  (100,) \rightarrow (1,100),
  \qquad
  (1,100) \text{ against } (100,1)
  \;\longrightarrow\; \mathbf{(100,100)} .
\]
So \texttt{y\_true - y\_pred} produces a $100\times100$ matrix of every
pairwise difference. Taking its mean square yields a number --- a perfectly
plausible-looking loss, computed over $10{,}000$ meaningless values instead of
$100$ real ones.

\emph{Nothing raises an error.} Training proceeds, the loss decreases, and the
model learns something unrelated to the task. The one-line defence:
\begin{center}
\texttt{assert y\_pred.shape == y\_true.shape}
\end{center}
This is the single most valuable assertion in applied machine learning.
\end{workedexample}

\begin{cautionbox}[Views, copies, and edits that vanish]
NumPy slicing returns a \emph{view}, not a copy: \texttt{b = a[0:3]} shares
memory with \texttt{a}, so writing to \texttt{b} changes \texttt{a}. Fancy
indexing (\texttt{a[[0,1,2]]}) returns a copy, so writing to it changes
nothing. The two look almost identical in code and behave oppositely, and
neither warns you.

In pandas the same ambiguity produces the notorious chained-assignment problem:
\texttt{df[df.x > 0]["y"] = 1} may modify a temporary and silently discard your
edit. Use \texttt{.loc} for any assignment, and call \texttt{.copy()} explicitly
when you intend a copy.
\end{cautionbox}

% ---- 3.7.1 NumPy arrays -------------------------------------------------------------------------
\subsubsection{NumPy arrays}

\begin{v3topic}{The NumPy ndarray}
    The \texttt{ndarray} is NumPy's central data structure: a contiguous block of typed memory with a shape, dtype, and strides\textsuperscript{\cite{fuhrerScientificComputingPython2021}}.
    All elements share the same \texttt{dtype}, enabling SIMD acceleration.
\visualseparator
    \begin{itemize}
        \item \texttt{np.array([1, 2, 3], dtype=np.float32)}
        \item \texttt{np.zeros((3, 4))}, \texttt{np.ones(10)}, \texttt{np.eye(4)}
        \item \texttt{np.arange(0, 1, 0.1)}, \texttt{np.linspace(0, 1, 100)}
        \item \texttt{arr.shape}, \texttt{arr.ndim}, \texttt{arr.dtype}, \texttt{arr.nbytes}
        \item \texttt{arr.reshape(2, -1)} --- reshape without copying data
    \end{itemize}

\end{v3topic}
\begin{v3topic}{Array dtypes}
    Choosing the right \texttt{dtype} balances precision against memory footprint, which matters at scale\textsuperscript{\cite{fuhrerScientificComputingPython2021}}.
\visualseparator
    \begin{tblr}{width=\linewidth,colspec={|Q[l,m,2.5cm]|Q[c,m]|X|}}
        \hline
        \textbf{dtype} & \textbf{Bytes} & \textbf{Use-case} \\
        \hline
        \texttt{float64} & 8 & Default; double precision \\
        \hline
        \texttt{float32} & 4 & Deep learning; halves memory \\
        \hline
        \texttt{int64}   & 8 & Large integer counts/indices \\
        \hline
        \texttt{int32}   & 4 & Compact integer arrays \\
        \hline
        \texttt{bool}    & 1 & Masks, logical operations \\
        \hline
    \end{tblr}
\captionof{table}{NumPy dtypes and their cost. Choosing \texttt{float32} over \texttt{float64} halves memory and is usually sufficient for training; it is not sufficient for accumulating sums (\cref{sec:math-numerics}).}
\label{tab:prog-numpy-dtypes}
\end{v3topic}

% ---- 3.7.2 Array indexing and slicing -----------------------------------------------------------
\subsubsection{Array indexing and slicing}

\begin{v3topic}{Indexing and Slicing}
    NumPy supports basic, fancy, and boolean indexing; slices return \emph{views} (no copy), fancy and boolean indexing return copies\textsuperscript{\cite{fuhrerScientificComputingPython2021}}.
\visualseparator
    \begin{tblr}{width=\linewidth,colspec={|Q[l,m,3.5cm]|X|}}
        \hline
        \textbf{Syntax} & \textbf{Meaning} \\
        \hline
        \texttt{a[2, 1]}            & Element at row 2, col 1 \\
        \hline
        \texttt{a[1:4, ::2]}        & Rows 1--3, every other column (view) \\
        \hline
        \texttt{a[[0, 2], :]}       & Rows 0 and 2 (fancy indexing, copy) \\
        \hline
        \texttt{a[a > 0]}           & All elements $> 0$ (boolean mask, copy) \\
        \hline
        \texttt{a[np.newaxis, :]}   & Add axis: shape $(n,) \to (1, n)$ \\
        \hline
        \texttt{a[..., 0]}          & Last axis = 0 for any number of leading axes \\
        \hline
    \end{tblr}
\captionof{table}{Indexing syntax. The distinction that bites: basic slicing returns a view sharing memory, fancy indexing returns a copy.}
\label{tab:prog-numpy-indexing}
\end{v3topic}

% ---- 3.7.3 Array operations ---------------------------------------------------------------------
\subsubsection{Array operations}

\begin{v3topic}{Vectorised Operations and ufuncs}
    NumPy operations are \emph{element-wise} by default; matrix products require \texttt{np.dot} or \texttt{@}.
    Reduction ufuncs (sum, max, etc.) accept an \texttt{axis} argument\textsuperscript{\cite{fuhrerScientificComputingPython2021}}.
\visualseparator
    \begin{tblr}{width=\linewidth,colspec={|Q[l,m,3.5cm]|X|}}
        \hline
        \textbf{Operation} & \textbf{Notes} \\
        \hline
        \texttt{a + b}, \texttt{a * b}       & Element-wise; broadcasting applies \\
        \hline
        \texttt{A @ B}                        & Matrix product (\texttt{np.matmul}) \\
        \hline
        \texttt{np.sum(a, axis=0)}            & Column-wise sum \\
        \hline
        \texttt{np.linalg.norm(v)}            & Euclidean norm \\
        \hline
        \texttt{np.sort(a, axis=-1)}          & Sort along last axis \\
        \hline
        \texttt{np.clip(a, 0, 1)}             & Clamp values to $[0, 1]$ \\
        \hline
    \end{tblr}
\captionof{table}{Vectorised operations. These are not merely shorthand for loops --- they dispatch to compiled routines, which is where the order-of-magnitude speedup comes from.}
\label{tab:prog-numpy-ufuncs}
\end{v3topic}

% ---- 3.7.4 Broadcasting -------------------------------------------------------------------------
\subsubsection{Broadcasting}

\begin{v3topic}{Broadcasting Rules}
    Broadcasting allows NumPy to operate on arrays with different shapes by implicitly expanding dimensions, without copying data\textsuperscript{\cite{fuhrerScientificComputingPython2021}}.
\visualseparator
    Two arrays are compatible if, for each trailing dimension, either the sizes match, or one of them is 1.
    \begin{equation}
        \textbf{A} \in \mathbb{R}^{m \times 1},\quad
        \textbf{b} \in \mathbb{R}^{1 \times n}
        \;\Longrightarrow\;
        \textbf{A} + \textbf{b} \in \mathbb{R}^{m \times n}
    \end{equation}
\where{$m$ and $n$ are the two surviving dimensions: a column of length $m$ and a
row of length $n$ broadcast to a full $m\times n$ grid. This is the mechanism
behind the silent shape bug described earlier in this chapter.}
    Example: \texttt{X - X.mean(axis=0)} subtracts a per-column mean from a matrix --- shape $(N, d) - (d,)$ is valid by broadcasting.
    A dimension of size 1 is \emph{stretched} (not copied) to match the other dimension.
\end{v3topic}

% ---- 3.8 Introduction to Pandas -----------------------------------------------------------------
\subsection{Introduction to Pandas}

% ---- 3.8.1 Series and DataFrames ----------------------------------------------------------------
\subsubsection{Series and DataFrames}

\begin{v3topic}{Series and DataFrames}
    Pandas provides two primary data structures: \texttt{Series} (1-D labelled array) and \texttt{DataFrame} (2-D labelled table)\textsuperscript{\cite{molinHandsonDataAnalysis2021}}.
    Both are backed by NumPy arrays and share an index.
\visualseparator
    \begin{itemize}
        \item \texttt{pd.Series([1, 2, 3], index=["a","b","c"])}
        \item \texttt{pd.DataFrame(\{"col1": [...], "col2": [...]\})}
        \item \texttt{df.shape}, \texttt{df.dtypes}, \texttt{df.head()}, \texttt{df.info()}
        \item \texttt{df.describe()} --- summary statistics
        \item \texttt{df["col"]} returns a Series; \texttt{df[["col1","col2"]]} returns a DataFrame
    \end{itemize}

\end{v3topic}
\begin{v3topic}{Reading and Writing Data}
    Pandas reads many formats natively; the resulting DataFrame is the entry point for all subsequent analysis\textsuperscript{\cite{molinHandsonDataAnalysis2021}}.
\visualseparator
    \begin{tblr}{width=\linewidth,colspec={|Q[l,m,3.5cm]|X|}}
        \hline
        \textbf{Function} & \textbf{Format} \\
        \hline
        \texttt{pd.read\_csv(path)}    & CSV / TSV \\
        \hline
        \texttt{pd.read\_parquet(path)} & Columnar Parquet \\
        \hline
        \texttt{pd.read\_excel(path)}  & Excel \texttt{.xlsx} \\
        \hline
        \texttt{pd.read\_sql(q, con)}  & SQL query result \\
        \hline
        \texttt{df.to\_csv(path)}      & Write CSV \\
        \hline
    \end{tblr}
\captionof{table}{Pandas I/O functions by format. Prefer a columnar binary format over CSV for anything re-read often: it preserves dtypes, which CSV does not.}
\label{tab:prog-pandas-io}
\end{v3topic}

% ---- 3.8.2 Data indexing and selection ----------------------------------------------------------
\subsubsection{Data indexing and selection}

\begin{v3topic}{loc, iloc, and Boolean Indexing}
    Pandas provides three main selection mechanisms\textsuperscript{\cite{molinHandsonDataAnalysis2021}}.
\visualseparator
    \begin{tblr}{width=\linewidth,colspec={|Q[l,m,2.5cm]|X[2]|X[2]|}}
        \hline
        \textbf{Accessor} & \textbf{Selects by} & \textbf{Example} \\
        \hline
        \texttt{df.loc}  & Label (index value or column name) & \texttt{df.loc["2020":"2022", "price"]} \\
        \hline
        \texttt{df.iloc} & Integer position & \texttt{df.iloc[0:5, 2:4]} \\
        \hline
        Boolean mask     & Condition evaluated element-wise & \texttt{df[df["age"] > 30]} \\
        \hline
        \texttt{df.query} & String expression (efficient) & \texttt{df.query("age > 30 and city == 'Berlin'")} \\
        \hline
    \end{tblr}
\captionof{table}{The three pandas accessors. Use \texttt{.loc} for any assignment --- chained indexing may write to a temporary and silently discard the edit.}
\label{tab:prog-pandas-accessors}
    Chained indexing (\texttt{df["a"]["b"]}) can cause \texttt{SettingWithCopyWarning}; prefer \texttt{.loc} for assignments.
\end{v3topic}

% ---- 3.8.3 Data cleaning and preparation --------------------------------------------------------
\subsubsection{Data cleaning and preparation}

\begin{v3topic}{Handling Missing Data}
    Pandas represents missing values as \texttt{NaN} (float) or \texttt{pd.NA}\textsuperscript{\cite{molinHandsonDataAnalysis2021}}.
\visualseparator
    \begin{itemize}
        \item \texttt{df.isna().sum()} --- count missing per column
        \item \texttt{df.dropna(subset=["col"])} --- drop rows with NaN in col
        \item \texttt{df.fillna(0)} --- replace NaN with constant
        \item \texttt{df["col"].fillna(df["col"].median())} --- impute with median
        \item \texttt{df.interpolate(method="time")} --- time-series interpolation
    \end{itemize}

\end{v3topic}
\begin{v3topic}{Type Conversion and Duplicates}
    Correct dtypes reduce memory and prevent silent bugs; duplicates inflate counts and bias models\textsuperscript{\cite{molinHandsonDataAnalysis2021}}.
\visualseparator
    \begin{itemize}
        \item \texttt{df["col"].astype("float32")} --- cast dtype
        \item \texttt{pd.to\_datetime(df["date"])} --- parse date strings
        \item \texttt{df["cat"].astype("category")} --- efficient categoricals
        \item \texttt{df.duplicated().sum()} --- count duplicate rows
        \item \texttt{df.drop\_duplicates(subset=["id"])} --- deduplicate on key
    \end{itemize}
\end{v3topic}

% ---- 3.8.4 Data aggregation and grouping --------------------------------------------------------
\subsubsection{Data aggregation and grouping}

\begin{v3topic}{GroupBy and Aggregation}
    The \texttt{groupby} / \texttt{agg} / \texttt{transform} pattern implements the \emph{split-apply-combine} strategy\textsuperscript{\cite{molinHandsonDataAnalysis2021}}.
\visualseparator
    \noindent\begin{tblr}{width=\linewidth,colspec={Q[l,m,5.3cm]X},
                           hlines,vlines}
        \hline
        \textbf{Pattern} & \textbf{Result} \\
        \hline
        \texttt{df.groupby("region").mean()} & Per-group means as new DataFrame \\
        \hline
        \texttt{df.groupby("region").\allowbreak
        agg(\{"sales": "sum",\allowbreak "profit": "mean"\})} & Multiple agg functions per column \\
        \hline
        \texttt{df.groupby("region")["sales"].\allowbreak
        transform("mean")} & Broadcast group mean back to original index \\
        \hline
        \texttt{df.pivot\_table(values="sales", index="region", columns="year", aggfunc="sum")} & Pivot for cross-tabulation \\
        \hline
    \end{tblr}
\captionof{table}{The split--apply--combine pattern. Most row-wise loops over a DataFrame are a group-by that has not been recognised yet.}
\label{tab:prog-pandas-groupby}
\end{v3topic}

% ---- 3.9 Data visualization with Pandas ---------------------------------------------------------
\subsection{Data visualization with Pandas}

% ---- 3.9.1 Plotting with Pandas -----------------------------------------------------------------
\subsubsection{Plotting with Pandas}

\begin{v3topic}{Plotting with Pandas}
    Pandas wraps Matplotlib's API via \texttt{DataFrame.plot}, enabling fast exploratory plots directly from a DataFrame without explicit Matplotlib calls\textsuperscript{\cite{molinHandsonDataAnalysis2021}}.
\visualseparator
    \begin{itemize}
        \item \texttt{df.plot(kind="line")} --- line chart
        \item \texttt{df["col"].plot(kind="hist", bins=30)} --- histogram
        \item \texttt{df.plot(kind="scatter", x="a", y="b")} --- scatter plot
        \item \texttt{df.plot(kind="box")} --- box plots per column
        \item \texttt{df.corr().style.background\_gradient()} --- annotated correlation heatmap
    \end{itemize}

\end{v3topic}
\begin{v3topic}{Customizing Plots}
    Every Pandas plot returns a Matplotlib \texttt{Axes} object, so all Matplotlib customisation applies\textsuperscript{\cite{molinHandsonDataAnalysis2021}}.
\visualseparator
    \begin{itemize}
        \item \texttt{ax.set\_title("...")}, \texttt{ax.set\_xlabel("...")}
        \item \texttt{ax.legend(loc="upper right")}
        \item \texttt{fig.tight\_layout()} --- prevent label clipping
        \item \texttt{fig.savefig("plot.pdf", dpi=300)} --- publication-quality export
        \item \texttt{plt.style.use("seaborn-v0\_8")} --- apply aesthetic theme
    \end{itemize}
\end{v3topic}

% ---- 3.10 Working with time series data ---------------------------------------------------------
\subsection{Working with time series data}

% ---- 3.10.1 Date and time handling --------------------------------------------------------------
\subsubsection{Date and time handling}

\begin{v3topic}{Date and Time in Pandas}
    Pandas provides a \texttt{DatetimeIndex} with nanosecond precision, enabling fast label-based time lookups and resampling\textsuperscript{\cite{molinHandsonDataAnalysis2021}}.
\visualseparator
    \begin{itemize}
        \item \texttt{pd.to\_datetime(series)} --- parse from strings
        \item \texttt{pd.date\_range("2020-01", periods=12, freq="M")} --- generate range
        \item \texttt{df.set\_index("date")} --- use datetime column as index
        \item \texttt{df.loc["2021":"2022"]} --- partial-string slicing on DatetimeIndex
        \item \texttt{ts.dt.year}, \texttt{ts.dt.dayofweek} --- accessor properties
    \end{itemize}

\end{v3topic}
\begin{v3topic}{Resampling and Frequency Conversion}
    \texttt{resample} aggregates time series to a coarser frequency; \texttt{asfreq} pads or forward-fills to a finer frequency\textsuperscript{\cite{molinHandsonDataAnalysis2021}}.
\visualseparator
    \begin{itemize}
        \item \texttt{df.resample("W").mean()} --- weekly mean
        \item \texttt{df.resample("Q").agg(\{"vol":"sum","price":"last"\})} --- quarterly
        \item \texttt{df.asfreq("D", method="ffill")} --- forward-fill daily
        \item \texttt{df.rolling(window=30).mean()} --- rolling 30-day average
        \item \texttt{df.ewm(span=10).mean()} --- exponentially weighted mean
    \end{itemize}
\end{v3topic}

% ---- 3.11 Advanced Pandas techniques ------------------------------------------------------------
\subsection{Advanced Pandas techniques}

% ---- 3.11.1 Merging and joining DataFrames ------------------------------------------------------
\subsubsection{Merging and joining DataFrames}

\begin{v3topic}{Merging and Joining DataFrames}
    Pandas mirrors SQL joins via \texttt{pd.merge}; \texttt{pd.concat} stacks DataFrames along an axis\textsuperscript{\cite{molinHandsonDataAnalysis2021}}.
\visualseparator
    \begin{tblr}{width=\linewidth,colspec={|Q[l,m,4cm]|X|}}
        \hline
        \textbf{Function / arg} & \textbf{Behaviour} \\
        \hline
        \texttt{pd.merge(A, B, on="id")}           & Inner join on column \texttt{id} \\
        \hline
        \texttt{pd.merge(A, B, how="left")}        & Left join; keeps all rows of A \\
        \hline
        \texttt{pd.merge(A, B, left\_on="a\_id", right\_on="b\_id")} & Join on differently-named keys \\
        \hline
        \texttt{pd.concat([df1, df2], axis=0)}     & Stack rows (append) \\
        \hline
        \texttt{pd.concat([df1, df2], axis=1)}     & Stack columns side-by-side \\
        \hline
        \texttt{df.join(other, on="key")}          & Join on index by default \\
        \hline
    \end{tblr}
\captionof{table}{Join behaviour. Always check row count before and after a merge: an unexpected increase means duplicate keys, and it is the most common silent data bug.}
\label{tab:prog-pandas-merge}
\end{v3topic}

% ---- 3.11.2 Pivot tables and cross-tabulations --------------------------------------------------
\subsubsection{Pivot tables and cross-tabulations}

\begin{v3topic}{Pivot Tables and Cross-Tabulations}
    Pivot tables reshape data from long (tidy) to wide format, aggregating values at the intersection of two categorical dimensions\textsuperscript{\cite{molinHandsonDataAnalysis2021}}.
\visualseparator
    \begin{tblr}{width=\linewidth,colspec={|Q[l,m,4.5cm]|X|}}
        \hline
        \textbf{Function} & \textbf{Notes} \\
        \hline
        \texttt{df.pivot\_table(values, index, columns, aggfunc)} & Generalised pivot with aggregation \\
        \hline
        \texttt{pd.crosstab(a, b)}                               & Frequency table of two categorical Series \\
        \hline
        \texttt{pd.crosstab(a, b, normalize="index")}            & Row-normalised proportions \\
        \hline
        \texttt{df.melt(id\_vars, value\_vars)}                   & Wide $\to$ long (tidy) format \\
        \hline
        \texttt{df.stack()} / \texttt{df.unstack()}              & Rotate innermost column/row level \\
        \hline
    \end{tblr}
\captionof{table}{Reshaping operations for moving between long and wide layouts.}
\label{tab:prog-pandas-pivot}
\end{v3topic}

% ---- 3.12 Performance optimization --------------------------------------------------------------
\subsection{Performance optimization}

% ---- 3.12.1 Efficient data manipulation ---------------------------------------------------------
\subsubsection{Efficient data manipulation}

\begin{v3topic}{Efficient Pandas Operations}
    Vectorised Pandas operations are typically $10\times$--$100\times$ faster than row-by-row Python loops\textsuperscript{\cite{molinHandsonDataAnalysis2021}}.
\visualseparator
    \begin{itemize}
        \item Prefer \texttt{df["col"].str.contains()} over looping
        \item \texttt{df.apply(fn, axis=1)} --- row-wise; use \texttt{np.where} or vectorised ops when possible
        \item \texttt{pd.cut} / \texttt{pd.qcut} --- vectorised binning
        \item Use categoricals for low-cardinality string columns: saves memory and speeds up groupby
        \item \texttt{df.pipe(fn)} --- chain multiple transformations cleanly
    \end{itemize}

\end{v3topic}
\begin{v3topic}{Memory Management}
    DataFrames can consume far more memory than their CSV source due to dtype widening\textsuperscript{\cite{molinHandsonDataAnalysis2021}}.
\visualseparator
    \begin{itemize}
        \item \texttt{df.memory\_usage(deep=True).sum()} --- total memory footprint
        \item Downcast integers: \texttt{pd.to\_numeric(s, downcast="integer")}
        \item Use \texttt{float32} instead of \texttt{float64} where precision allows
        \item Read large files in chunks: \texttt{pd.read\_csv(path, chunksize=10\_000)}
        \item \texttt{del df} and \texttt{gc.collect()} --- release memory explicitly
    \end{itemize}
\end{v3topic}

% ==================================================================================================
\section{SciKit-Learn}

\begin{v3topic}{Scikit-learn: Estimator API}
    Scikit-learn implements a uniform \emph{Estimator API}: every model exposes \texttt{fit}, \texttt{predict} (or \texttt{transform}), and \texttt{score}\textsuperscript{\cite{raschkaPythonMachineLearning2015}}.
    This uniformity enables drop-in model swapping and automatic pipeline composition.
\visualseparator
    \begin{itemize}
        \item \texttt{model.fit(X\_train, y\_train)} --- learn parameters
        \item \texttt{model.predict(X\_test)} --- apply learned model
        \item \texttt{model.score(X\_test, y\_test)} --- default metric (accuracy / $R^2$)
        \item \texttt{transformer.fit\_transform(X)} --- fit and apply in one step
        \item Hyperparameter tuning via \texttt{GridSearchCV} / \texttt{RandomizedSearchCV}
    \end{itemize}

\end{v3topic}
\begin{v3topic}{Scikit-learn: Key Modules}
    Scikit-learn organises its functionality into sub-modules that cover the full ML workflow\textsuperscript{\cite{raschkaPythonMachineLearning2015}}.
\visualseparator
    \begin{tblr}{width=\linewidth,colspec={|Q[l,m,3.5cm]|X|}}
        \hline
        \textbf{Module} & \textbf{Contents} \\
        \hline
        \texttt{sklearn.linear\_model}   & LR, Ridge, Lasso, SGD \\
        \hline
        \texttt{sklearn.tree}            & Decision trees \\
        \hline
        \texttt{sklearn.ensemble}        & Random forest, gradient boosting \\
        \hline
        \texttt{sklearn.svm}             & Support vector machines \\
        \hline
        \texttt{sklearn.preprocessing}   & StandardScaler, OneHotEncoder, LabelEncoder \\
        \hline
        \texttt{sklearn.pipeline}        & \texttt{Pipeline}, \texttt{ColumnTransformer} \\
        \hline
        \texttt{sklearn.model\_selection} & Cross-validation, search \\
        \hline
        \texttt{sklearn.metrics}         & Accuracy, F1, ROC-AUC, MSE, \ldots \\
        \hline
    \end{tblr}
\captionof{table}{The scikit-learn module map. Everything follows the same fit/transform/predict interface, which is what makes pipelines composable.}
\label{tab:prog-sklearn-modules}

\end{v3topic}
\begin{v3topic}{Scikit-learn Pipeline}
    A \texttt{Pipeline} chains preprocessing and model steps; \texttt{fit} on training data, \texttt{predict} on test data --- no data leakage\textsuperscript{\cite{raschkaPythonMachineLearning2015}}.
\visualseparator
    \begin{center}
    \begin{tikzpicture}[
        node distance=0.3cm,
        every node/.style={draw, rounded corners, minimum height=0.65cm, minimum width=1.8cm, font=\sffamily\footnotesize, align=center},
        arr/.style={-{Stealth[length=2mm]}, thick}
    ]
        \node (imp)   {Imputer};
        \node (scl)   [right=of imp]  {Scaler};
        \node (enc)   [right=of scl]  {Encoder};
        \node (feat)  [right=of enc]  {Feature\\Selection};
        \node (model) [right=of feat] {Estimator};
        \draw[arr] (imp) -- (scl);
        \draw[arr] (scl) -- (enc);
        \draw[arr] (enc) -- (feat);
        \draw[arr] (feat) -- (model);
    \end{tikzpicture}
\captionof{figure}{A scikit-learn pipeline. Its real value is not tidiness:
because every transform is fitted inside the pipeline, cross-validation refits
them on each training fold, which structurally prevents the preprocessing leak
described in \cref{ch:data-analytics}.}
\label{fig:prog-sklearn-pipeline}
    \end{center}
    \texttt{ColumnTransformer} applies different transformers to different feature subsets in parallel before feeding results to the final estimator.
\end{v3topic}

% ==================================================================================================
\section{PyTorch}

\begin{v3topic}{PyTorch Tensors}
    PyTorch's core data structure is the \texttt{Tensor}: a multi-dimensional array that lives on CPU or GPU and records operations for automatic differentiation\textsuperscript{\cite{raschkaPythonMachineLearning2015}}.
\visualseparator
    \begin{itemize}
        \item \texttt{torch.tensor([1.0, 2.0], dtype=torch.float32)}
        \item \texttt{t.to("cuda")} --- move to GPU; \texttt{t.cpu()} --- move back
        \item \texttt{t.shape}, \texttt{t.dtype}, \texttt{t.device}
        \item \texttt{t.numpy()} --- zero-copy view as NumPy array (CPU only)
        \item \texttt{torch.zeros}, \texttt{torch.randn}, \texttt{torch.arange}
    \end{itemize}

\end{v3topic}
\begin{v3topic}{Autograd and the Computation Graph}
    Setting \texttt{requires\_grad=True} on a tensor tells PyTorch to track all operations on it, building a \emph{computation graph} from which gradients are derived via reverse-mode automatic differentiation\textsuperscript{\cite{raschkaPythonMachineLearning2015}}.
\visualseparator
    \begin{equation}
        \frac{\partial \mathcal{L}}{\partial \mathbf{w}} = \text{backprop through graph}
    \end{equation}
\where{$\mathcal{L}$ is the loss and $\mathbf{w}$ the parameters; the framework
records the operations performed in the forward pass and replays them in
reverse to obtain this derivative (\cref{sec:dl-backprop}).}
    \begin{itemize}
        \item \texttt{loss.backward()} --- compute all gradients
        \item \texttt{w.grad} --- accumulated gradient $\partial\mathcal{L}/\partial w$
        \item \texttt{optimizer.zero\_grad()} --- reset before each step
        \item \texttt{torch.no\_grad():} --- disable graph for inference
    \end{itemize}

\end{v3topic}
\begin{v3topic}{nn.Module and the Training Loop}
    All neural network models subclass \texttt{torch.nn.Module}, implementing \texttt{forward()}\textsuperscript{\cite{raschkaPythonMachineLearning2015}}.
    The standard training loop is explicit: forward pass $\to$ loss $\to$ backward pass $\to$ optimizer step.
\visualseparator
    \begin{tblr}{width=\linewidth,colspec={|Q[l,m,3.5cm]|X|}}
        \hline
        \textbf{Step} & \textbf{Code} \\
        \hline
        Define model    & \texttt{class Net(nn.Module): def forward(self, x): ...} \\
        \hline
        Loss function   & \texttt{criterion = nn.CrossEntropyLoss()} \\
        \hline
        Optimizer       & \texttt{optimizer = torch.optim.Adam(model.parameters(), lr=1e-3)} \\
        \hline
        Forward pass    & \texttt{output = model(x\_batch)} \\
        \hline
        Compute loss    & \texttt{loss = criterion(output, y\_batch)} \\
        \hline
        Backward pass   & \texttt{optimizer.zero\_grad(); loss.backward()} \\
        \hline
        Update weights  & \texttt{optimizer.step()} \\
        \hline
    \end{tblr}
\captionof{table}{The PyTorch training loop, step by step. Forgetting \texttt{zero\_grad()} accumulates gradients across batches --- a bug that trains, slowly and wrongly.}
\label{tab:prog-pytorch-loop}
    \texttt{torch.utils.data.DataLoader} wraps any \texttt{Dataset} and handles batching, shuffling, and parallel loading.
\end{v3topic}

% ==================================================================================================
%
\section{Deep Dive: Reproducible ML Software}
\begin{v3topic}{Notebook to Package}
Keep exploration in notebooks, but move reusable transformations, models, and
metrics into typed modules with documented interfaces.
A command-line entry point plus a versioned configuration makes the experiment
rerunnable without executing cells in a remembered order.

\end{v3topic}
\begin{v3topic}{Test Properties, Not Only Examples}
Example tests check known cases.
Property and metamorphic tests check invariants: shuffling rows should not
change a permutation-invariant aggregate; serialization should preserve a
prediction within tolerance; a monotone transform should preserve ranks.
These tests catch families of failures that one golden input misses.

\end{v3topic}
\begin{v3topic}{Operational Skeleton}
Build a minimal repository with locked dependencies, configuration validation,
structured logging, data-schema checks, unit tests, one end-to-end smoke test,
an environment capture, and a single command that trains and evaluates a
baseline.
Record code revision, data identifier, seed, hardware, and output artifact in
the same run manifest.
\end{v3topic}


\begin{workedexample}[The record that makes a result reproducible]
An experiment is reproducible when someone else, months later, can obtain your
number. That requires recording seven things, and the omission of any one of
them is enough to break it:

\smallskip
\begin{tblr}{
  width=\linewidth,
  colspec={Q[l,0.22\linewidth] X[l]},
  row{1}={font=\bfseries,bg=AlmanachMist},
  hlines,vlines}
Record & Because without it \\
Data version or hash & ``the dataset'' silently changes under you \\
Split indices or seed & a different split is a different experiment \\
All random seeds & Python, NumPy, and the framework each have their own \\
Library versions & a minor release can change a default and the result \\
Hardware and precision & GPU reductions are non-deterministic in ordering \\
Full configuration & a hyperparameter you did not log is one you cannot
  restore \\
Exact metric definition & ``accuracy'' hides averaging and threshold choices \\
\end{tblr}
\smallskip

\emph{Note the honest limit.} Even with all seven, bitwise reproducibility on a
GPU is not guaranteed --- floating-point addition is not associative, and
parallel reductions may sum in a different order between runs
(\cref{sec:math-numerics}). What you can promise is that the result is
reproducible \emph{within run-to-run variance}, which is why you report a
spread over seeds rather than one number.
\end{workedexample}

\begin{cautionbox}[The notebook that cannot be re-run]
A Jupyter notebook records the order cells were \emph{written}, not the order
they were \emph{executed}. A notebook that produces a beautiful result may
depend on a cell that was edited after it ran, a variable defined in a cell
since deleted, or an execution order nobody remembers. Before trusting any
notebook result: restart the kernel and run all cells top to bottom. If it does
not reproduce, the result does not exist yet.
\end{cautionbox}

\begin{practicallab}[Make one experiment genuinely reproducible]
\textbf{Goal.} Turn a notebook that works into an experiment that survives.

\textbf{Protocol.}
\begin{enumerate}
  \item Take any small training script you have. Run it twice without changing
        anything. Record whether the two numbers are identical.
  \item Set every seed --- Python's \texttt{random}, NumPy, and the
        framework's. Run twice more. Are they identical now? On GPU, probably
        still not; note that.
  \item Run five seeds and report mean and standard deviation. This spread is
        the resolution of your experiment: any improvement smaller than it is
        not yet demonstrated.
  \item Write out a JSON record with the seven fields from the worked example.
  \item Move the data loading and metric into a module with a test. Write one
        property test: shuffling the rows must not change the metric.
  \item Restart and run everything from a clean shell. Delete every cached
        artefact first.
\end{enumerate}

\textbf{Deliverable.} The JSON record, the seed-spread number, and one passing
property test. \textbf{The point:} step 3 tells you how large a difference has
to be before you are allowed to call it an improvement.
\end{practicallab}

\begin{checkpoint}
A colleague reports that their model improved from $0.882$ to $0.889$ accuracy.
Using this chapter, list three questions you would ask before believing the
improvement is real, and name the single number they should have reported
alongside it.
\end{checkpoint}

\chapterreview{Specify}{Implement}{Test}{Maintain}

\section{Further Reading}
\begin{itemize}
    \item Lutz, M. (2017). Learning Python (5th ed.). O'Reilly.
    \item Mathes, E. (2019). Python Crash Course (2nd ed.). No Starch Press.
\end{itemize}
